% Figure: a semicircular curved cantilever of radius R, fixed at one end and
% carrying a vertical load P at the free end. The angle phi locates a generic
% cross-section; the moment arm of P about that section is the horizontal
% distance from the load to the section. Used in the Castigliano worked
% example on a curved member.
\begin{figure}[htbp]
\centering
\begin{tikzpicture}[
  member/.style={draw=engblue, very thick},
  arm/.style={draw=enggray, dashed, thin},
  loadarrow/.style={draw=engred, -{Stealth}, thick},
  wall/.style={draw=engblack, thick},
  label/.style={font=\small},
]
% Semicircle of radius R = 2.4 units, centre at O = (0,0).
% Fixed end at the top, angle phi measured from the top (fixed) end.
% Parametrise a point on the arc by angle a from positive x-axis:
% take the arc from (0, R) [top] sweeping clockwise to (0, -R) [bottom free].
% Use a in degrees: point = (R sin a, R cos a) for a in 0..180 gives
% top (0,R) at a=0 down to bottom (0,-R) at a=180, passing right side.
\def\R{2.4}

% the arc
\draw[member] plot[domain=0:180, samples=80]
  ({\R*sin(\x)}, {\R*cos(\x)});

% centre O
\fill[engblack] (0,0) circle (1.2pt);
\node[font=\scriptsize, anchor=east] at (-0.1,0) {$O$};

% fixed (built-in) end at top (0,R): hatched wall
\draw[wall] (-0.7,\R) -- (0.7,\R);
\foreach \i in {-0.6,-0.4,-0.2,0,0.2,0.4,0.6} {
  \draw[wall, thin] (\i,\R) -- (\i-0.18,\R+0.28);
}
\node[label, anchor=south] at (0,\R+0.35) {fixed end};

% free end at bottom (0,-R), with load P downward
\fill[engblue] (0,-\R) circle (1.6pt);
\draw[loadarrow] (0,-\R) -- (0,-\R-1.0);
\node[engred, label, anchor=west] at (0.1,-\R-0.6) {$P$};
\node[label, anchor=east] at (-0.15,-\R+0.05) {free end};

% a generic section at angle phi (pick a = 60 deg from top)
\def\a{60}
\coordinate (S) at ({\R*sin(\a)}, {\R*cos(\a)});
\fill[engamber] (S) circle (1.6pt);
% radius line O==S to indicate the angle
\draw[arm] (0,0) -- (S);
\node[engamber, font=\scriptsize, anchor=south west] at (S) {section $\phi$};

% horizontal moment arm of P about the section S:
% the load acts at x = 0, the section is at x = R sin(phi); the moment arm
% is the horizontal distance R sin(phi).
\draw[arm] (S) -- ({\R*sin(\a)}, -\R);
\node[enggray, font=\scriptsize, anchor=west] at ({\R*sin(\a)+0.1}, -\R+0.5)
  {arm $=R\sin\phi$};

\end{tikzpicture}
\caption[Semicircular curved cantilever under a tip load]{A semicircular
curved member of radius $R$, built in at the top and free at the bottom,
carrying a vertical load $P$ at its free end. The angle $\phi$ measured from
the fixed end locates a generic cross-section. The bending moment at that
section is the load times its horizontal moment arm, $M(\phi)=PR\sin\phi$,
which the curved-member form of Castigliano's theorem integrates around the
arc to give the tip deflection. The curvature of the member, rather than its
slenderness, sets the form of the moment, and the energy integral runs over
arc length $R\,d\phi$ rather than a straight coordinate.}
\label{fig:vol03ch05:curved-member}
\end{figure}
